\nuevaReceta{Boquerones en Vinagre}{30 min (+ reposo y congelación)}{receta:boquerones_vinagre}

    \begin{minipage}[t]{0.35\textwidth}
        \section*{Ingredientes}
        \begin{itemize}[leftmargin=*]
            \item 500 g de boquerones frescos
            \item 250 ml de vinagre de vino blanco
            \item 100 ml de agua fría
            \item 3 Dientes de ajo
            \item Perejil fresco
            \item Aceite de oliva virgen extra
            \item Sal
            \item Hielo
        \end{itemize}
    \end{minipage}
    \hfill
    \begin{minipage}[t]{0.6\textwidth}
        \section*{Preparación}
        \begin{enumerate}
            \item \textbf{Limpiar los boquerones:} Retirar la cabeza, las vísceras y la espina central. Separarlos en dos lomos y lavarlos con agua fría.
            \item \textbf{Desangrar:} Colocar los lomos en un recipiente con agua, hielo y una pizca de sal. Dejarlos entre 30 y 60 minutos en la nevera, cambiando el agua si fuera necesario, hasta que queden limpios.
            \item \textbf{Congelar:} Escurrirlos, secarlos y congelarlos durante al menos 5 días en un congelador doméstico adecuado. Descongelarlos después lentamente dentro de la nevera.
            \item \textbf{Preparar el marinado:} Mezclar el vinagre con el agua fría y una cucharadita de sal.
            \item \textbf{Marinar:} Colocar los lomos con la piel hacia abajo, cubrirlos por completo con la mezcla y dejarlos en la nevera durante 4-6 horas, hasta que estén blancos y firmes.
            \item \textbf{Escurrir:} Retirar el marinado y secar ligeramente los boquerones para que el aliño no quede aguado.
            \item \textbf{Aliñar:} Disponerlos por capas con el ajo y el perejil muy picados. Cubrirlos con aceite de oliva y dejarlos reposar en la nevera al menos 2 horas antes de servir.
        \end{enumerate}
    \end{minipage}

    \vspace{1em}
    \begin{tcolorbox}[colback=orange!5, colframe=orange!50, title=Importante, size=small]
        La congelación preventiva es necesaria para reducir el riesgo de anisakis. Conserva siempre los boquerones refrigerados, cubiertos de aceite y en un recipiente limpio y cerrado.
    \end{tcolorbox}
\end{receta}
